\section{Classificação}%
\subsection{Término}%
\begin{frame}[label=main]
	\begin{itemize}
		\item Alguns autores restringem a definição de algoritmo para 
		procedimentos que terminam (eventualmente).
			\subitem{Um algoritmo pode repetir um procedimento infinitamente}
		\pause\vfill
		\item Se o tamanho de um procedimento não é conhecido, não é possível 
		determinar se ele terminará (Marvin Minsky).
			\subitem{Provado por Alan Turing em 1936}
		\pause\vfill
		\item Para algoritmos intermináveis o sucesso não pode ser determinado 
		pela interpretação da resposta e sim por condições impostas pelo próprio
		desenvolvedor do algoritmo durante sua execução.
	\end{itemize}
\end{frame}

\subsection{Representação}%
\begin{frame}[label=main]
	\begin{columns}
		\column{.65\textwidth}
			\begin{itemize}[<+->]
				\item Linguagem Natural
					\subitem{Se estiver \alert<1-1>{quente}, tome \alert<1>{remédio}}
				\item Pseudo-código
					\subitem{\textbf{se} \emph{temp $> 37^{\circ}{\rm C}$} \textbf{então} \emph{tomar(aspirina)}}
				\item Diagrama
				\item Tabela de Controle
					\vfill
					\begin{center}
						\small
						\begin{tabular}{|c|c|c|c|}
							\hline
							\textbf{ASCII} & \textbf{dec} & \textbf{hex} & \textbf{vetor} \\\hline
							A & 65 & 41 & abacate \\\hline
							B & 66 & 42 & banana \\\hline
							$\cdots$ & $\cdots$ & $\cdots$ & $\cdots$ \\\hline
							P & 80 & 50 & pêra \\\hline
						\end{tabular}
					\end{center}
					\vfill
				\item Linguagem de Programação
					\subitem{\textbf{if}(\emph{temp$>$37}) \textbf{take}(aspirin)}
			\end{itemize}
		\column{.3\textwidth}
			\includegraphics<5->[width=\textwidth]{flowchart}%
			%source: http://en.wikipedia.org/wiki/File:FlowchartExample.png
	\end{columns}
\end{frame}

\subsection{Implementação}%
\begin{frame}[label=main]
	\begin{itemize}
	    \item Iterativo%: estruturas de repetição (laços, pilhas, etc.)
	    \item Recursivo%: invoca a \alert<1>{si mesmo} até que certa condição seja satisfeita
	    \vfill\pause
		\item Lógico%: deduções lógicas controladas
	    \vfill\pause
		\item Serial%: cada instrução é executada em \alert<3>{seqüência}
		\item Paralela%: várias instruções executadas \alert<3>{ao mesmo tempo}
	    \vfill\pause
		\item Determinístico%: decisão \alert<4>{exata} a cada passo
		\item Probabilístico%: decisão \alert<4>{provável} em passo(s)
	    \vfill\pause
		\item Exato%: resposta \alert<5>{exata}
		\item Aproximado%: resposta \alert<5>{próxima} a verdadeira solução
	\end{itemize}
\end{frame}

\begin{frame}[fragile,label=graduacao]
	\framesubtitle{Recursividade}
	
	Lida com uma função que pode invocar a si própria
	\vspace{-.7\baselineskip}
	\begin{columns}[T]
		\column{.4\textwidth}
			\begin{lstlisting}[title=Fatorial iterativo]
				fatorial(n)
				    resultado := 1
				    enquanto n > 1
				      resultado := n*resultado
				      n :=  n - 1
				    retorna resultado
			\end{lstlisting}
			\pause
		\column{.4\textwidth}
			\begin{lstlisting}[title=Fatorial recursivo]
				fatorial(n)
				    se n < 2
				      retorna 1
				    senao
				      retorna n*fatorial(n-1)
			\end{lstlisting}
	\end{columns}
	\vfill\pause
	Recursão vs Iteração?
	% \begin{itemize}
	%	\item a implementação iterativa tende a ser ligeiramente mais rápida na 
	%	prática do que a implementação recursiva%, uma vez que uma implementação recursiva precisa registrar o estado atual do processamento de maneira que ela possa continuar de onde parou após a conclusão de cada nova excecução subordinada do procedimento recursivo. Esta ação consome tempo e memória. (Note que a implementação de uma função fatorial para números naturais pequenos é mais rápida quando se usa uma tabela de busca.)
	%	\item existem problemas cujas soluções são inerentemente recursivas%, já que elas precisam manter registros de estados anteriores. Um exemplo é o percurso de uma árvore; outros exemplos incluem a função de Ackermann e algoritmos de divisão e conquista tais como o Quicksort. Todos estes algoritmos podem ser implementados iterativamente com a ajuda de uma pilha, mas o uso de uma pilha, de certa forma, anula as vantagens das soluções iterativas.
	%	\item nas linguagens de programação modernas o espaço disponível para o 
	%	fluxo de controle é geralmente bem menor que o espaço da pilha%, e algoritmos recursivos tendem a necessitar de mais espaço na pilha do que algoritmos iterativos.
	% \end{itemize}
\end{frame}

\subsection{Paradigma}%
\begin{frame}[label=main]
	\begin{itemize}[<+->]
		\item Força Bruta%: tentar todas as soluções possíveis
		\item Divisão e conquista%: reduzir repetidamente o problema, até que este seja pequeno o suficiente para ser resolvido
		\item Programação dinâmica%: reutilizar o cálculo de solução já resolvidas
		\item Algoritmo ganancioso%: escolher a solução que parece ser mais adequada no momento
		\item Programação linear%: otimização de uma função linear
		\item Redução%: transformação em outro problema de mais fácil solução
		\item Busca e enumeração%: exploração de possíveis soluções
	\end{itemize}
\end{frame}

\subsection{Complexidade}
\begin{frame}[label=main]
	Porque analisar um algoritmo?
	\begin{enumerate}
		\item Para saber quais os recursos necessários
		\item Para avaliar seu desempenho (e comparar diferentes algoritmos)
	\end{enumerate}
	\pause\vfill
	Mas analisar o que?
	\begin{description}
		\item[tempo] de execução
		\item[espaço] uso de memória
	\end{description}
	\pause\vfill
	\begin{alertblock}{Importante!}
		Análise de algoritmos compara \alert{algoritmos} e não \emph{programas}.
	\end{alertblock}
\end{frame}

\begin{frame}[label=graduacao]
	Análise de um algoritmo particular
	\begin{itemize}
		\item Qual é o custo de usar um dado algoritmo para resolver um problema específico?
		\item Características que devem ser investigadas:
			\begin{itemize}
				\item análise do número de vezes que cada parte do algoritmo 
				deve ser executada $\rightarrow$ \alert<1>{tempo}
				\item estudo da quantidade de memória necessária $\rightarrow$ \alert<1>{espaço}
			\end{itemize}
	\end{itemize}
	\vfill\pause
	Análise de uma classe de algoritmos
	\begin{itemize}
		\item Qual é o algoritmo de menor custo possível para resolver um problema particular?
		\item Toda uma família de algoritmos é investigada (busca pelo melhor possível)
		\item Coloca-se limites para a complexidade computacional dos algoritmos pertencentes à classe
	\end{itemize}
\end{frame}

\begin{frame}[label=graduacao]
	Custo de um algoritmo
	\begin{itemize}
		\item ao determinar o menor custo possível para resolver problemas de uma 
		classe, tem-se a medida da dificuldade inerente para resolver o problema.
		\item quando o custo de um algoritmo é igual ao menor custo possível, o 
		algoritmo é considerado ótimo (para a medida de custo considerada)
	\end{itemize}
	\vfill\pause
	Podem existir vários algoritmos para resolver o mesmo problema
	\begin{itemize}
		\item se a mesma medida de custo é aplicada a diferentes algoritmos, 
		então é possível compará-los e escolher o mais adequado
		\item se há vários algoritmos, todos com um custo de execução dentro da 
		mesma ordem de grandeza, pode-se considerar tanto os custos reais das 
		operações como os custos não aparentes
			\vspace{-1\baselineskip}
			\begin{columns}[T]
				\column{.45\textwidth}
					\begin{itemize}
						\setlength{\itemsep}{0pt}
						\item alocação de memória
						\item indexação
					\end{itemize}
				\column{.45\textwidth}
					\begin{itemize}
						\setlength{\itemsep}{0pt}
						\item carga
						\item etc.
					\end{itemize}
			\end{columns}
	\end{itemize}
\end{frame}

\begin{frame}[label=graduacao]
	Análise pela execução
		\begin{itemize}
			\item a eficiência do programa depende da linguagem (compilada ou interpretada)
			\item depende do sistema operacional
			\item depende do hardware (quantidade de memória, velocidade do 
			processador e arquitetura)
		\end{itemize}
	\vfill\pause
	\begin{beamerboxesrounded}{}
		Útil nas comparações entre programas em máquinas específicas pois, além 
		dos custos do algoritmo, são comparados os custos não aparentes (alocação 
		de memória, indexação, carga de arquivos, etc)
	\end{beamerboxesrounded}
\end{frame}

\begin{frame}[label=main]
	Modelo matemático: ``contar instruções'' 
		\begin{itemize}
			\item Não depende do computador nem da implementação
			\item Determinado em função do custo das operações realizadas e do
			tamanho da entrada
		\end{itemize}
	\pause \vfill
	Caracteriza de forma simples o desempenho do algoritmo e permite comparações 
	entre algoritmos distintos que resolvem um mesmo problema.
\end{frame}

\begin{frame}[label=main]
	Exemplo: descobrir o maior número em uma lista
	\\\vfill\pause
	Função de custo: \alert{atribuição}
		\subitem{hipótese: todas as atribuição custam o mesmo}
	\vfill\pause
	\begin{columns}[T]
		\column{.42\textwidth}
			\underline{Algoritmo} \\
			lista $\gets$ \textbf{entrada}()\\
			maior $\gets -\infty$\\
			\textbf{para} i $\in [1,n]$ \\
			\hspace{1em}maior $\gets$ \textbf{max}(maior, lista[i])
			\pause
		\column{.30\textwidth}
			\underline{Tempo} \\
			$T\cdot 1$ \\
			$T\cdot 1$ \\
			$T\cdot n$ \\
			$T\cdot n$\\
			\pause
			$\overline{g_T(n)=2n+2}$
		\pause
		\column{.25\textwidth}
			\underline{Espaço} \\
			$S{\cdot}n$ \\
			$S{\cdot}1$ \\
			$S{\cdot}2$ \\
			$S{\cdot}0$ \\
			\pause
			$\overline{g_S(n)=n+3}$
	\end{columns}
\end{frame}

\begin{frame}[fragile,label=main]
	Exemplo: Qual o custo do algoritmo de soma?
		\begin{center}
			\begin{minipage}{.3\textwidth}
				\begin{lstlisting}
					soma := 0
					para i /in [1, n]
					    soma := soma + i
				\end{lstlisting}
			\end{minipage}
		\end{center}
	\pause
	Considerando somente as \alert<2>{atribuições} como operações relevantes.
	\\\vfill\pause
	\begin{columns}[T]
		\column{.45\textwidth}
			Custo (tempo): \alert{$g(n) = 2n + 1$}
			\begin{itemize}
				\item[1] na inicialização
				\item[2n] no laço
			\end{itemize}
		\pause
		\column{.45\textwidth}
			Custo (espaço): \alert{$g(n) = 3$}
			\subitem[3]{variáveis ($soma$, $i$, $n$)}
	\end{columns}
\end{frame}

\begin{frame}[fragile,label=main]
	Exemplo: Qual o custo do algoritmo de divisão de elementos de uma matriz 
	quadrada por uma constante $k$?
	\begin{center}
	\begin{minipage}{.4\textwidth}
	\begin{lstlisting}
		para i /in [1, n]
		    para j /in [1, n]
		        A[i][j] := A[i][j]/k;
	\end{lstlisting}
	\end{minipage}
	\end{center}
	\pause
	Considerando somente as \alert<2>{atribuições} como operações relevantes
	\\\vfill\pause
	\begin{columns}[T]
		\column{.51\textwidth}
			Custo (tempo): \alert{$g(n)=2n^2+n$}  
				\begin{itemize}
					\item[$n$] no laço externo ($i$)
					\item[$n^2$] no laço interno ($j$)
					\item[$n^2$] no laço interno ($A_{ij}$)
				\end{itemize}
		\pause
		\column{.44\textwidth}
			Custo (espaço): \alert{$g(n) = n^2+4$} 
				\begin{itemize}
					\item[4] variáveis simples ($i$, $n$, $j$, $k$)
					\item[$n^2$] matriz ($A$)
				\end{itemize}
	\end{columns}
\end{frame}

\newcommand{\custo}[1]{$g(#1) = 2{\cdot}#1^2 + #1$}%
\begin{frame}[label=main]
	Considerando a função \custo{n}, o crescimento dos termos em função de $n$:
	\begin{itemize}
		\item \custo{1}   	$ = 3$\pause
		\item \custo{5}   	$ = 55$\pause
		\item \custo{10}  	$ = 210$
		\item \custo{100} 	$ = 20100$
		\item \custo{1000}	$ = 2001000$
		\item $\cdots$
	\end{itemize}
	\vfill\pause
	Para valores de $n$ muito grandes, o \alert{primeiro} termo é \alert{dominante}...
\end{frame}